\documentclass[12pt]{article}
\usepackage{amsmath, amssymb}
\usepackage[margin=1in]{geometry}
\setlength{\parskip}{6pt}
\title{QUBO Formulation: Benchmark Job Selection Problem}
\date{}
\begin{document}
\maketitle

\subsection*{Benchmark Job Selection Problem}

\paragraph{Problem Statement.}
Given a set of $N$ candidate benchmark jobs, each characterized by a coverage value $v_i$, a compute time $t_i$, and pairwise redundancy penalties $p_{ij}$, select a subset of jobs to maximize the total benchmark coverage value while penalizing redundant coverage, subject to a total compute-time budget $T$.

\paragraph{Decision Variables.}
For each candidate job $i \in \{0, \dots, N-1\}$, introduce a binary variable $x_i \in \{0,1\}$, where $x_i = 1$ if job $i$ is selected, and $x_i = 0$ otherwise.

\paragraph{QUBO Derivation.}
\textbf{Step 1.} The original objective is a maximization problem. To cast it as a QUBO minimization, we negate the objective function:
\begin{align}
    \min_{x \in \{0,1\}^N} \left( -\sum_{i=0}^{N-1} v_i x_i + \sum_{0 \leq i < j \leq N-1} p_{ij} x_i x_j \right).
\end{align}

\textbf{Step 2.} The compute-time budget is expressed as a linear inequality constraint:
\begin{align}
    \sum_{i=0}^{N-1} t_i x_i \leq T.
\end{align}

\textbf{Step 3.} We enforce the constraint via a quadratic penalty term with weight $\lambda > 0$:
\begin{align}
    P(x) = \lambda \left( \sum_{i=0}^{N-1} t_i x_i - T \right)^2.
\end{align}
Expanding the square and applying the idempotent property $x_i^2 = x_i$ for binary variables yields:
\begin{align}
    P(x) &= \lambda \left( \sum_{i=0}^{N-1} t_i^2 x_i^2 + \sum_{i \neq j} t_i t_j x_i x_j - 2T \sum_{i=0}^{N-1} t_i x_i + T^2 \right) \\
    &= \lambda \left( \sum_{i=0}^{N-1} (t_i^2 - 2T t_i) x_i + 2 \sum_{0 \leq i < j \leq N-1} t_i t_j x_i x_j + T^2 \right).
\end{align}
The general penalty contribution to the Hamiltonian is therefore:
\begin{align}
    P(x) = \sum_{i=0}^{N-1} \lambda (t_i^2 - 2T t_i) x_i + 2\lambda \sum_{0 \leq i < j \leq N-1} t_i t_j x_i x_j + \lambda T^2.
\end{align}

\textbf{Step 4.} Combining the negated objective and the penalty term gives the full QUBO Hamiltonian $H(x)$:
\begin{align}
    H(x) = \sum_{i=0}^{N-1} \left( -v_i + \lambda(t_i^2 - 2T t_i) \right) x_i + \sum_{0 \leq i < j \leq N-1} \left( p_{ij} + 2\lambda t_i t_j \right) x_i x_j + \lambda T^2.
\end{align}

\textbf{Step 5.} Reading off the upper-triangular QUBO matrix entries $Q_{i,j}$ and the constant offset $c$ from $H(x) = \sum_{i \leq j} Q_{i,j} x_i x_j + c$, we obtain the closed-form expressions:
\begin{align}
    Q_{i,i} &= -v_i + \lambda(t_i^2 - 2T t_i), \quad i \in \{0, \dots, N-1\}, \\
    Q_{i,j} &= p_{ij} + 2\lambda t_i t_j, \quad 0 \leq i < j \leq N-1, \\
    Q_{i,j} &= 0, \quad i > j, \\
    c &= \lambda T^2.
\end{align}

\paragraph{Penalty Weight Justification.}
The penalty weight $\lambda$ must be chosen large enough so that the cost of violating the budget constraint strictly exceeds any possible gain in the benchmark coverage objective. Let $V_{\max} = \max_i v_i$. A sufficient condition is $\lambda > V_{\max}$ (or more conservatively $\lambda > \sum_i v_i$). This ensures that any assignment violating $\sum_i t_i x_i \leq T$ incurs a penalty larger than the maximum achievable objective value, thereby forcing the global minimum of $H(x)$ to lie within the feasible region.

\paragraph{Correctness Argument.}
Minimizing $H(x)$ over $\{0,1\}^N$ recovers the optimal solution to the original problem. For any feasible assignment satisfying $\sum_i t_i x_i \leq T$, the penalty term $P(x)$ evaluates to zero, and $H(x)$ reduces exactly to the negated original objective. Thus, minimizing $H(x)$ over feasible points is equivalent to maximizing the original objective. For any infeasible assignment, $P(x)$ is strictly positive. By choosing $\lambda$ according to the justification above, the energy of every infeasible point exceeds the energy of the optimal feasible point. Consequently, the global minimum of $H(x)$ must correspond to a feasible assignment that maximizes the original objective function.

\paragraph{Illustrative Example.}
Consider an instance with $N=3$ jobs. We instantiate the parameters as follows: coverage values $v_0=10, v_1=15, v_2=12$; redundancy penalties $p_{01}=2, p_{02}=3, p_{12}=2$; compute times $t_0=3, t_1=4, t_2=5$; budget $T=8$; and penalty weight $\lambda=100$. Substituting these values into the general closed-form expressions yields the concrete QUBO coefficients:
\begin{align}
    Q_{0,0} &= -10 + 100(3^2 - 2\cdot8\cdot3) = -3910, \\
    Q_{1,1} &= -15 + 100(4^2 - 2\cdot8\cdot4) = -4815, \\
    Q_{2,2} &= -12 + 100(5^2 - 2\cdot8\cdot5) = -5512, \\
    Q_{0,1} &= 2 + 2\cdot100\cdot3\cdot4 = 2402, \\
    Q_{0,2} &= 3 + 2\cdot100\cdot3\cdot5 = 3003, \\
    Q_{1,2} &= 2 + 2\cdot100\cdot4\cdot5 = 4002, \\
    c &= 100 \cdot 8^2 = 6400.
\end{align}
The resulting Hamiltonian is:
\begin{align}
    H(x) = -3910 x_0 - 4815 x_1 - 5512 x_2 + 2402 x_0 x_1 + 3003 x_0 x_2 + 4002 x_1 x_2 + 6400.
\end{align}
Evaluating $H(x)$ over all feasible assignments (those with $\sum_i t_i x_i \leq 8$) gives:
\begin{align}
    H(0,0,0) &= 6400, \quad H(1,0,0) = 2490, \quad H(0,1,0) = 1585, \quad H(0,0,1) = 888, \\
    H(1,1,0) &= -923, \quad H(1,0,1) = -19.
\end{align}
The assignment $x^* = (1,1,0)$ yields the minimum energy $H(x^*) = -923$, corresponding to selecting jobs $0$ and $1$. This matches the optimal benchmark coverage value of $v_0 + v_1 - p_{01} = 10 + 15 - 2 = 23$, confirming the formulation's correctness on this instance.

\end{document}